\documentclass[11pt]{article}

\usepackage[margin=1in]{geometry}
\usepackage{setspace}
\usepackage{graphicx}
\usepackage{booktabs}
\usepackage{amsmath}
\usepackage{url}
\usepackage[hidelinks]{hyperref}
\usepackage{natbib}
\usepackage{microtype}

\singlespacing

\title{\textbf{Remote Work and the Changing Geography of U.S. Housing Demand}}
\author{
Sai Krishna Varshith Danda \and
Gayathri Kota \and
Eshitha Chowdary Kaka \and
Zain Khan
}
\date{\today}

\begin{document}

\maketitle

% ============================================================
\section{Introduction}
% ============================================================

% PURPOSE:
% Establish the broad relationship between employment location,
% commuting, and residential location choice, then narrow toward
% the COVID-19 remote-work shock and your specific research problem.
%
% Possible flow:
%
% 1. Residential location and commuting have historically been linked.
%
% 2. COVID-19 caused a sudden increase in working from home.
%
% 3. Remote work reduced the importance of commuting for many workers,
%    potentially changing where households were willing to live.
%
% 4. Research from 2020--2022 documented changes in urban housing
%    demand, including movement away from dense urban cores.
%
% 5. However, the persistence of these changes after the acute
%    pandemic period is less clear.
%
% 6. Introduce the focus of this review:
%    - urban/suburban housing price gradients
%    - remote work as a mechanism
%    - neighborhood-level differences
%    - persistence into the post-pandemic period
%
% End by explaining how the review informs the project's research
% questions, datasets, and methodology.


% ============================================================
\section{Remote Work and Residential Location}
% ============================================================

% Establish the theoretical/mechanistic connection between
% remote work and housing demand.
%
% Topics to investigate:
%
% - How commuting costs affect residential location decisions.
% - Why remote work could increase demand for locations farther
%   from employment centers.
% - Changes in the prevalence of remote work before and after COVID.
% - Which occupations can be performed remotely.
% - Differences in remote-work capability across workers,
%   occupations, and income groups.
%
% Important starting sources from proposal:
%
% Barrero, Bloom, and Davis -- Why Working From Home Will Stick
%
% Dingel and Neiman -- How Many Jobs Can Be Done at Home?


% ============================================================
\section{The Pandemic and the Geography of Housing Demand}
% ============================================================

% Review evidence that COVID-era housing demand changed
% geographically.
%
% Do NOT organize this as:
% "Ramani found X. Gupta found Y. Ahrend found Z."
%
% Instead, synthesize the evidence around common findings
% and methodological differences.


\subsection{The ``Donut Effect''}

% Ramani and Bloom are central here.
%
% Discuss:
% - Definition of the Donut Effect
% - Movement away from dense central areas
% - Housing price/rent changes
% - Migration/change-of-address evidence
% - Differences between urban cores and suburbs
%
% Compare their evidence with other studies rather than
% treating the paper in isolation.


\subsection{Changes in Urban Housing Price Gradients}

% Connect directly to Sub-Question 1.
%
% Discuss literature measuring:
% - distance from CBD
% - housing price gradients
% - rent gradients
% - density
% - changes before vs. during COVID
%
% Gupta et al. is an important starting point.
%
% Questions:
% Did prices farther from CBDs grow faster?
% Did the traditional urban price gradient flatten?
% Were effects similar for rents and home values?


\subsection{Evidence Across Cities and Countries}

% Compare findings across different geographic contexts.
%
% OECD paper can be important here.
%
% Investigate whether:
% - effects occur outside the largest U.S. cities
% - magnitude differs across metropolitan areas
% - Sun Belt cities behave differently
% - housing supply affects the response
%
% This subsection can help justify your focus on
% high-growth Sun Belt metropolitan areas.


% ============================================================
\section{Remote Work as a Mechanism}
% ============================================================

% Connect directly to Sub-Question 2.
%
% The literature may establish that suburban prices changed,
% but that does not automatically establish remote work as
% the mechanism.
%
% Synthesize research connecting:
%
% remote-work prevalence
%          |
%          v
% residential mobility
%          |
%          v
% neighborhood housing demand
%          |
%          v
% prices/rents
%
% Discuss alternative explanations where supported by literature:
% - fiscal/monetary conditions
% - housing supply constraints
% - household preferences for space
% - migration
% - demographic changes
%
% IMPORTANT:
% Distinguish causal findings from associations.
%
% Your proposal already recognizes that remote-work share and
% housing appreciation can be bidirectionally related.


\subsection{Measuring Remote-Work Exposure}

% Review ways previous researchers measure remote work:
%
% - ACS "worked from home" share
% - occupation-based remote-work feasibility
% - surveys such as SWAA
% - employer/workplace data if found
%
% Compare advantages and limitations.
%
% This section should help justify the variable(s) ultimately
% selected for your analysis.


% ============================================================
\section{Persistence After the Initial Pandemic Shock}
% ============================================================

% THIS IS ESPECIALLY IMPORTANT FOR YOUR PROJECT.
%
% Earlier studies heavily emphasize 2020--2022.
% Your project extends analysis through 2025.
%
% Review newer literature examining:
% - persistence of remote work
% - return-to-office
% - migration after 2022
% - housing prices after 2022
% - whether urban cores recovered
%
% Ramani, Alcedo, and Bloom (2024) is an important starting point.
%
% Compare evidence supporting:
%
% 1. persistence
% 2. partial reversal
% 3. geographic heterogeneity
%
% Avoid forcing the literature into one conclusion if evidence
% is mixed.


% ============================================================
\section{Neighborhood Heterogeneity}
% ============================================================

% Connect to Sub-Question 3.
%
% Investigate literature showing that neighborhoods did not
% experience the pandemic housing market uniformly.
%
% Characteristics potentially associated with different outcomes:
%
% - distance from CBD
% - density
% - income
% - remote-work share
% - occupation
% - housing type
% - housing supply
% - new construction
%
% Look for studies using:
% - ZIP codes
% - census tracts
% - neighborhoods
% - clustering or other unsupervised methods
%
% This section should establish why analyzing heterogeneous
% trajectories is useful rather than relying only on an
% average metropolitan effect.


% ============================================================
\section{Data and Methods Used in Prior Research}
% ============================================================

% This section directly supports the eventual Methods section
% of the capstone.


\subsection{Housing Market Data}

% Compare datasets used by prior researchers.
%
% Zillow:
% - ZHVI
% - ZORI
%
% Other datasets you encounter:
% - transaction data
% - Redfin
% - CoreLogic
% - Census
% etc.
%
% Explain strengths/limitations relevant to your project.


\subsection{Demographic and Remote-Work Data}

% ACS
% SWAA
% occupation-based measures
% migration data
%
% Discuss geographic resolution and temporal availability.


\subsection{Empirical Methods}

% Review methods relevant to your analysis:
%
% - panel regression
% - fixed effects
% - event studies
% - spatial analysis
% - price-gradient estimation
% - clustering
%
% Focus on WHY researchers use these methods and what
% assumptions/limitations they have.


% ============================================================
\section{Gaps and Open Questions}
% ============================================================

% Synthesize what remains unresolved.
%
% This section should naturally motivate your project without
% making unsupported claims of novelty.
%
% Questions to investigate from the literature:
%
% - How persistent was the Donut Effect after 2022?
%
% - Did return-to-office substantially reverse the spatial
%   reallocation?
%
% - Are effects different in high-growth Sun Belt metros?
%
% - How strongly does neighborhood remote-work prevalence
%   predict subsequent housing-price trajectories?
%
% - Are average CBD-distance effects hiding distinct groups
%   of neighborhood trajectories?
%
% - How important are housing supply differences?
%
% These should lead naturally into your project's three
% research questions.


% ============================================================
\section{Conclusions}
% ============================================================

% Synthesize rather than simply repeat earlier sections.
%
% Possible structure:
%
% 1. What the literature establishes relatively strongly.
%
% 2. What remains uncertain or contested.
%
% 3. What datasets and methods appear most appropriate.
%
% 4. How these gaps motivate examining 2018--2025 ZIP-level
%    housing trajectories in Sun Belt metropolitan areas.
%
% The conclusion should make the connection between the
% literature and your eventual empirical analysis clear.


% ============================================================
% References
% ============================================================

\newpage
\bibliographystyle{apalike}
\bibliography{references}

\end{document}